% ============================================================
%  PCEOT402 – Electrical Machines
%  Module 3: Synchronous Machines – Complete Study Notes
%  Built from Parts 1, 2 & 3 and styled to Anthropic branding
% ============================================================
\documentclass[12pt,a4paper]{article}

% ── packages ────────────────────────────────────────────────
\usepackage[a4paper,
            top=2.5cm, bottom=2.5cm,
            left=2.2cm, right=2.2cm]{geometry}
\usepackage{amsmath,amssymb,mathtools}
\usepackage{booktabs,multirow,array,longtable}
\usepackage{xcolor}
\usepackage{tcolorbox}
\tcbuselibrary{breakable,theorems}
\usepackage{titlesec}
\usepackage{enumitem}
\usepackage{fancyhdr}
\usepackage{hyperref}
\usepackage{microtype}
\usepackage{parskip}
\usepackage{graphicx}
\usepackage{tikz}
\usetikzlibrary{arrows.meta,positioning,calc,decorations.pathreplacing}
\usepackage{circuitikz}
\usepackage{caption}
\usepackage{subcaption}
\usepackage{tabularx}
\usepackage{colortbl}
\usepackage{float}
\usepackage{fontspec}
\usepackage{setspace}
\usepackage{pdfpages}
\usepackage{pdflscape}

% ── Anthropic Brand Colour Palette ───────────────────────────
\definecolor{anthDark}{HTML}{141413}
\definecolor{anthLight}{HTML}{FAF9F5}
\definecolor{anthMidGray}{HTML}{B0AEA5}
\definecolor{anthLightGray}{HTML}{E8E6DC}
\definecolor{anthOrange}{HTML}{D97757}
\definecolor{anthBlue}{HTML}{6A9BCC}
\definecolor{anthGreen}{HTML}{788C5D}
\definecolor{anthOrangeLight}{HTML}{F5DDD5}
\definecolor{anthBluePale}{HTML}{DDEAF5}
\definecolor{anthGreenPale}{HTML}{DDE5D2}
\definecolor{anthDarkLight}{HTML}{3A3A38}

% ── typography ───────────────────────────────────────────────
\IfFileExists{assets/fonts/Lora/LoraVar.ttf}{
  \setmainfont{LoraVar}[
    Path=assets/fonts/Lora/,
    Extension=.ttf,
    UprightFont=*,
    ItalicFont=LoraItalicVar,
    BoldFont=*,
    BoldItalicFont=LoraItalicVar
  ]
}{
  \setmainfont{Liberation Serif}
}
\IfFontExistsTF{Poppins}{
  \newfontfamily\HeadingFont{Poppins}
}{
  \newfontfamily\HeadingFont{Liberation Sans}
}

% ── page look ────────────────────────────────────────────────
\pagecolor{anthLight}
\color{anthDark}
\setstretch{1.08}

% ── hyperref setup ───────────────────────────────────────────
\hypersetup{
  colorlinks=true,
  linkcolor=anthDark,
  urlcolor=anthBlue,
  pdftitle={PCEOT402 – Module 3: Synchronous Machines},
  pdfauthor={Electrical Machines},
}

% ── section/subsection formatting ────────────────────────────
\titleformat{\section}[block]
  {\HeadingFont\color{anthLight}\Large\bfseries}
  {\thesection\quad}
  {0pt}
  {\colorbox{anthDark}\bgroup\parbox{\dimexpr\linewidth-2\fboxsep\relax}}
  [\egroup\vspace{2pt}]

\titleformat{\subsection}
  {\HeadingFont\color{anthDark}\large\bfseries}
  {\thesubsection}
  {0.6em}
  {}

\titleformat{\subsubsection}
  {\HeadingFont\color{anthBlue}\normalsize\bfseries}
  {\thesubsubsection}
  {0.5em}
  {}

\titlespacing*{\section}{0pt}{16pt}{8pt}
\titlespacing*{\subsection}{0pt}{12pt}{4pt}
\titlespacing*{\subsubsection}{0pt}{8pt}{3pt}

% ── tcolorbox environments ───────────────────────────────────
\tcbset{
  studybox/.style={
    breakable,
    colback=anthLightGray,
    colframe=anthDark,
    colbacktitle=anthDark,
    coltitle=anthLight,
    fonttitle=\HeadingFont\bfseries,
    arc=3pt,
    boxrule=0.8pt,
    left=8pt,
    right=8pt,
    top=8pt,
    bottom=8pt,
  },
  formulabox/.style={
    breakable,
    colback=anthBluePale,
    colframe=anthBlue,
    arc=3pt,
    boxrule=0.7pt,
    left=10pt,
    right=10pt,
    top=6pt,
    bottom=6pt,
  },
  definebox/.style={
    breakable,
    colback=anthDark!6,
    colframe=anthDark,
    colbacktitle=anthBluePale,
    coltitle=anthDark,
    fonttitle=\HeadingFont\bfseries\small,
    arc=3pt,
    boxrule=0.6pt,
    left=8pt,
    right=8pt,
    top=6pt,
    bottom=6pt,
  },
  importantbox/.style={
    breakable,
    colback=anthOrangeLight,
    colframe=anthOrange,
    colbacktitle=anthOrangeLight,
    coltitle=anthOrange,
    fonttitle=\HeadingFont\bfseries\small,
    arc=3pt,
    boxrule=0.7pt,
    left=8pt,
    right=8pt,
    top=8pt,
    bottom=8pt,
  },
  notebox/.style={
    breakable,
    colback=anthGreenPale,
    colframe=anthGreen,
    arc=2pt,
    boxrule=0.5pt,
    left=6pt,
    right=6pt,
    top=4pt,
    bottom=4pt,
  },
  examplebox/.style={
    breakable,
    colback=anthLightGray,
    colframe=anthMidGray,
    colbacktitle=anthMidGray,
    coltitle=anthDark,
    fonttitle=\HeadingFont\bfseries,
    arc=3pt,
    boxrule=0.6pt,
    left=8pt,
    right=8pt,
    top=8pt,
    bottom=8pt,
  },
}

\newtcolorbox{studybox}[1][]{studybox, title={#1}}
\newtcolorbox{formulabox}{formulabox}
\newtcolorbox[auto counter,number within=section]{definebox}[1][]{
  definebox, title={Definition~\thetcbcounter\ — #1}}
\newtcolorbox{importantbox}[1][Key Point]{importantbox, title={#1}}
\newtcolorbox{notebox}{notebox}
\newtcolorbox{examplebox}[1][]{examplebox, title={#1}}

% ── table helpers ────────────────────────────────────────────
\newcolumntype{C}[1]{>{\centering\arraybackslash}p{#1}}
\newcolumntype{L}[1]{>{\raggedright\arraybackslash}p{#1}}
\newcolumntype{R}[1]{>{\raggedleft\arraybackslash}p{#1}}

\newcommand{\thead}[1]{\cellcolor{anthDark}\color{anthLight}\HeadingFont\bfseries #1}
\newcommand{\trow}[0]{\rowcolor{anthLightGray}}
\newcommand{\troww}[0]{\rowcolor{anthBluePale}}

% ── header / footer ──────────────────────────────────────────
\pagestyle{fancy}
\fancyhf{}
\fancyhead[L]{\small\color{anthDark}\HeadingFont\textbf{PCEOT402 – Electrical Machines}}
\fancyhead[R]{\small\color{anthDark}\HeadingFont Module 3: Synchronous Machines}
\fancyfoot[C]{\color{anthMidGray}\small\thepage}
\setlength{\headheight}{15pt}
\renewcommand{\headrulewidth}{1pt}
\renewcommand{\headrule}{\color{anthOrange}\hrule width\headwidth height\headrulewidth}

% ── misc macros ──────────────────────────────────────────────
\newcommand{\degr}{^\circ}
\newcommand{\Eph}{E_{\mathrm{ph}}}
\newcommand{\Vt}{V_t}
\newcommand{\Ef}{E_f}
\newcommand{\Ia}{I_a}
\newcommand{\Ra}{R_a}
\newcommand{\Xs}{X_s}
\newcommand{\Phim}{\Phi_m}
\newcommand{\kp}{k_p}
\newcommand{\kd}{k_d}
\newcommand{\kw}{k_w}

% ============================================================
\begin{document}

% ── TITLE PAGE ──────────────────────────────────────────────
\thispagestyle{empty}
\begin{tikzpicture}[remember picture, overlay]
  \fill[anthDark] (current page.north west) rectangle
    ([yshift=-7.8cm]current page.north east);
  \fill[anthOrange!70] ([yshift=-7.8cm]current page.north west) rectangle
    ([yshift=-8.25cm]current page.north east);
  \fill[anthBlue!22] ([xshift=-1.2cm,yshift=-1.2cm]current page.north east) circle (1.9cm);
  \fill[anthGreen!20] ([xshift=1.7cm,yshift=1.4cm]current page.south west) circle (2.3cm);
\end{tikzpicture}

\vspace*{0.7cm}
\begin{center}
  {\color{anthLight}\HeadingFont\fontsize{28}{34}\selectfont\bfseries PCEOT402}\\[6pt]
  {\color{anthLight}\HeadingFont\fontsize{22}{28}\selectfont\bfseries Electrical Machines}\\[10pt]
  {\color{anthOrange!92}\HeadingFont\fontsize{16}{22}\selectfont Module 3 — Synchronous Machines}\\[4pt]
  {\color{anthLightGray}\large Complete Study Notes from Parts 1, 2 \& 3}
\end{center}

\vspace{4.4cm}

\begin{tcolorbox}[studybox={Topics Covered in This Module}]
\begin{enumerate}[leftmargin=*, label=\textcolor{anthOrange}{\HeadingFont\bfseries\arabic*.}]
  \item Principle of operation of the 3-phase synchronous generator
  \item Classification and constructional features of alternators
  \item Speed of rotation, pole relation and generated frequency
  \item Pitch factor, distribution factor, winding factor and EMF equation
  \item Open-circuit characteristics and no-load operation
  \item Armature reaction and synchronous reactance / impedance
  \item Equivalent circuit, load phasor diagrams and voltage regulation
  \item Predetermination of regulation by EMF and MMF methods
  \item Parallel operation of alternators and synchronising conditions
  \item Dark-lamp, bright-lamp and synchroscope methods
  \item Synchronising current, synchronising torque and synchronising power
  \item Synchronous motor construction and principle of operation
  \item Why the motor is not self-starting and practical starting methods
  \item Equivalent circuit of the synchronous motor and effect of excitation
  \item Original slide appendix for exact source-page reference
\end{enumerate}
\end{tcolorbox}

\vspace{0.8cm}
\begin{notebox}
\textbf{How this document is organised:} the main body is written as clean study notes in the
style of your Module 2 document. An appendix at the end embeds the original slide decks so
that no source material is lost.
\end{notebox}

\newpage
\setcounter{page}{1}
\tableofcontents
\newpage

% ============================================================
\section{Principle of Operation of a 3-Phase Synchronous Generator}
% ============================================================

\begin{definebox}[Synchronous Generator]
A \textbf{synchronous generator} or \textbf{alternator} is an AC generator in which the rotor
magnetic field is excited by DC and driven by a prime mover so that a 3-phase emf is induced in
the stator armature windings.
\end{definebox}

\subsection{Physical Idea}

\begin{itemize}[leftmargin=*]
  \item The rotor carries the \textbf{field winding}; the stator carries the \textbf{armature winding}.
  \item DC excitation establishes fixed rotor poles.
  \item When the rotor is driven mechanically, the rotor field sweeps past the stationary stator conductors.
  \item Hence the flux linking each stator phase changes continuously and an emf is induced according to Faraday's law.
  \item Since the stator contains a 3-phase distributed winding, the induced voltages are 3-phase and displaced by $120\degr$ electrical.
\end{itemize}

\subsection{Faraday-Law Basis}

\begin{formulabox}
\[
  e = N\frac{d\Phi}{dt}
\]
For a rotating coil in a sinusoidally distributed field, the flux linkage varies sinusoidally; hence the induced emf is also sinusoidal.
\end{formulabox}

\begin{importantbox}[AC Generation in One Line]
Mechanical rotation \(\rightarrow\) changing magnetic flux linkage \(\rightarrow\) induced alternating emf.
\end{importantbox}

\subsection{Why Slip Rings are Used}

\begin{itemize}[leftmargin=*]
  \item If the output were taken directly from the rotating coil, the connecting wires would twist.
  \item Hence \textbf{slip rings} and \textbf{brushes} are used for a rotating coil arrangement.
  \item In a practical alternator, however, the armature is stationary and only the rotor field needs DC excitation through slip rings.
\end{itemize}

% ============================================================
\section{Classification and Construction of Alternators}
% ============================================================

\subsection{Classification}

\begin{center}
\begin{tabular}{L{3.4cm} L{8.8cm}}
\toprule
\thead{Basis} & \thead{Types}\\
\midrule
\trow Prime mover & Hydro generators, turbo generators, diesel-engine generators, gas-turbine generators\\
Number of phases & Single-phase, three-phase\\
\trow Rotating member & Rotating armature, rotating field\\
Rotor construction & Salient-pole machine, non-salient (cylindrical / turbo) machine\\
\bottomrule
\end{tabular}
\end{center}

\subsection{Standard Construction}

\begin{itemize}[leftmargin=*]
  \item \textbf{Stator:} houses the armature conductors in slots on the inner periphery.
  \item \textbf{Rotor:} carries the field poles and field winding supplied with DC.
  \item \textbf{Excitation:} supplied either from an external DC source through slip rings or from a shaft-mounted exciter.
\end{itemize}

\begin{importantbox}[Why Rotating Field and Stationary Armature are Preferred]
\begin{itemize}[leftmargin=*]
  \item High voltage can be generated at the stationary armature terminals.
  \item Better insulation is possible for the armature winding.
  \item Collection of output current is easier.
  \item Rotor becomes lighter, especially for large machines.
  \item Only two slip rings are required for the field supply.
\end{itemize}
\end{importantbox}

\subsection{Salient-Pole vs Cylindrical Rotor}

\begin{center}
\begin{tabular}{L{3.0cm} L{4.9cm} L{4.9cm}}
\toprule
\thead{Feature} & \thead{Salient-Pole Rotor} & \thead{Cylindrical Rotor}\\
\midrule
\trow Rotor form & Projecting poles mounted on rotor body & Smooth solid cylindrical rotor\\
Air gap & Non-uniform & Uniform\\
\trow Shape & Large diameter, short axial length & Small diameter, long axial length\\
Speed range & Low / medium speed & High speed\\
\trow Number of poles & Usually many poles & Usually 2 or 4 poles\\
Mechanical strength & Lower & Very high\\
\trow Typical use & Hydroelectric stations & Steam / turbo alternators\\
Cooling & Natural / air cooling & Hydrogen or forced cooling possible\\
\bottomrule
\end{tabular}
\end{center}

\subsection{Constructional Notes from the Slides}

\begin{itemize}[leftmargin=*]
  \item The armature core is built from laminated silicon-steel sheets to reduce hysteresis and eddy-current losses.
  \item Cylindrical rotors have low windage loss, good balancing and quieter operation.
  \item Salient-pole machines are well suited to low-speed prime movers such as water turbines.
\end{itemize}

% ============================================================
\section{Speed of Rotation, Frequency and Winding Factors}
% ============================================================

\subsection{Speed-Frequency Relation}

\begin{formulabox}
\[
  f = \frac{P N_s}{120}
  \qquad\Longleftrightarrow\qquad
  N_s = \frac{120f}{P}
\]
where \(f\) is electrical frequency in hertz, \(N_s\) is synchronous speed in rpm, and \(P\) is the number of poles.
\end{formulabox}

\begin{notebox}
The generator is called \textbf{synchronous} because the electrical frequency is locked to rotor speed and pole number.
\end{notebox}

\subsection{Pitch Factor, Distribution Factor and Winding Factor}

\begin{definebox}[Pitch Factor / Coil-Span Factor]
The \textbf{pitch factor} \(\kp\) accounts for the reduction in generated emf when a coil is short-pitched.
\end{definebox}

\begin{formulabox}
\[
  \kp = \cos\frac{\alpha}{2}
\]
where \(\alpha\) is the chording angle.
\end{formulabox}

\begin{definebox}[Distribution Factor]
The \textbf{distribution factor} \(\kd\) accounts for the reduction in resultant emf when the phase winding is distributed over several slots.
\end{definebox}

\begin{formulabox}
\[
  \kd = \frac{\sin\left(\dfrac{m\beta}{2}\right)}{m\sin\left(\dfrac{\beta}{2}\right)}
\]
where \(m\) is the number of slots per pole per phase and \(\beta\) is the slot angle.
\end{formulabox}

\begin{importantbox}[Winding Factor]
\[
  \kw = \kp \kd
\]
For sinusoidal flux distribution, \(\kw\) is used directly in the emf equation.
\end{importantbox}

\subsection{EMF Equation}

\begin{formulabox}
\[
  \Eph = 4.44 \, \kw \, f \, \Phim \, T
\]
where \(T\) is the number of series turns per phase and \(\Phim\) is flux per pole in weber.
\end{formulabox}

\[
  \Eph = 4.44 \, \kp \kd \, f \, \Phim \, T
\]

\begin{examplebox}[Example — Generated EMF per Phase]
\textbf{Given:} \(f = 50\,\mathrm{Hz}\), \(\Phim = 0.05\,\mathrm{Wb}\), \(T = 100\) turns/phase,
\(\kp = 0.96\), \(\kd = 0.95\). Find the generated emf per phase.

\[
  \kw = \kp\kd = 0.96 \times 0.95 = 0.912
\]
\[
  \Eph = 4.44 \times 0.912 \times 50 \times 0.05 \times 100
  \approx 1012\ \mathrm{V}
\]
If the winding is star connected, line voltage \(\approx \sqrt{3}\,\Eph \approx 1753\ \mathrm{V}\).
\end{examplebox}

% ============================================================
\section{Open-Circuit Characteristics and No-Load Operation}
% ============================================================

\begin{definebox}[Open-Circuit Characteristic]
The \textbf{open-circuit characteristic} (OCC) of an alternator is the graph between the no-load generated emf and field current when the machine runs at constant rated speed with the armature open-circuited.
\end{definebox}

\subsection{Features of OCC}

\begin{itemize}[leftmargin=*]
  \item At zero field current, a small residual emf is present because of residual magnetism.
  \item Initially, the OCC is almost linear; this region is called the \textbf{air-gap line}.
  \item At higher excitation, magnetic saturation sets in and the curve bends.
\end{itemize}

\subsection{Why OCC is Important}

\begin{itemize}[leftmargin=*]
  \item It gives the relation between field current and generated voltage.
  \item It is used for predetermination of voltage regulation.
  \item Together with short-circuit data, it helps determine synchronous impedance.
\end{itemize}

\begin{notebox}
On no-load, the terminal voltage of the alternator is practically equal to the internally generated emf per phase because armature current is zero.
\end{notebox}

% ============================================================
\section{Armature Reaction, Leakage Reactance and Synchronous Impedance}
% ============================================================

\begin{definebox}[Armature Reaction]
\textbf{Armature reaction} is the effect of armature-flux mmf on the main field flux in the alternator.
\end{definebox}

\subsection{Nature of Armature Reaction}

\begin{center}
\begin{tabular}{L{4.5cm} L{7.5cm}}
\toprule
\thead{Load Condition} & \thead{Effect of Armature Reaction}\\
\midrule
\trow Unity power factor & Cross-magnetising\\
Zero power factor lagging & Demagnetising\\
\trow Zero power factor leading & Magnetising\\
General lagging power factor & Partly demagnetising and partly cross-magnetising\\
\trow General leading power factor & Partly magnetising and partly cross-magnetising\\
\bottomrule
\end{tabular}
\end{center}

\subsection{Leakage Reactance and Synchronous Reactance}

\begin{itemize}[leftmargin=*]
  \item Some flux produced by armature current does not cross the air-gap in the main path; it forms \textbf{leakage flux}.
  \item Leakage flux produces a voltage drop represented by the \textbf{leakage reactance} \(X_l\).
  \item Armature reaction is often represented by an equivalent reactance \(X_a\).
\end{itemize}

\begin{formulabox}
\[
  X_s = X_l + X_a
  \qquad\text{(synchronous reactance)}
\]
\[
  Z_s = \sqrt{\Ra^2 + X_s^2}
  \qquad\text{(synchronous impedance)}
\]
\end{formulabox}

\begin{importantbox}[Practical Meaning]
The synchronous reactance combines the effect of armature leakage and armature reaction into one equivalent series reactance.
\end{importantbox}

% ============================================================
\section{Equivalent Circuit, Phasor Diagram and Load Characteristics}
% ============================================================

\subsection{Per-Phase Equivalent Circuit of an Alternator}

\begin{center}
\begin{circuitikz}[american voltages, scale=0.95]
  \draw
  (0,0) to[sV, l=$\Ef$] (0,3)
  (0,3) to[R, l=$\Ra$] (3,3)
        to[L, l=$jX_s$] (6,3)
        to[short, -o] (7,3)
  (0,0) to[short, -o] (7,0);
  \draw (7,1.5) node[right] {$\Vt$};
  \draw[-{Latex[length=2mm]}] (4,3.45) -- (5.4,3.45) node[midway, above] {$\Ia$};
\end{circuitikz}
\end{center}

\begin{formulabox}
\[
  \Ef = \Vt + \Ia(\Ra + jX_s)
\]
If \(\Ra\) is neglected,
\[
  \Ef \approx \Vt + j\Ia X_s
\]
\end{formulabox}

\subsection{Phasor Behaviour with Different Loads}

\begin{itemize}[leftmargin=*]
  \item \textbf{Lagging load:} armature current lags terminal voltage; terminal voltage tends to fall because the reactive drop adds strongly.
  \item \textbf{Unity power factor load:} armature current is nearly in phase with terminal voltage; armature reaction is mainly cross-magnetising.
  \item \textbf{Leading load:} armature reaction becomes magnetising; terminal voltage may rise with load and voltage regulation can become negative.
\end{itemize}

\subsection{Load Characteristic}

\begin{itemize}[leftmargin=*]
  \item The load characteristic is the relation between terminal voltage and armature current at constant speed and constant excitation.
  \item For lagging power-factor load, terminal voltage falls more rapidly.
  \item For leading power-factor load, voltage fall is smaller and may even reverse.
\end{itemize}

% ============================================================
\section{Voltage Regulation and Predetermination Methods}
% ============================================================

\begin{definebox}[Voltage Regulation]
Voltage regulation of an alternator is the change in terminal voltage when the full-load is removed, while speed and excitation are kept constant, expressed as a fraction or percentage of full-load voltage.
\end{definebox}

\begin{formulabox}
\[
  \%\,\text{Voltage Regulation}
  =
  \frac{E_0 - \Vt}{\Vt}\times 100
\]
where \(E_0\) is the no-load emf corresponding to the same excitation and \(\Vt\) is the full-load terminal voltage.
\end{formulabox}

\subsection{EMF Method (Synchronous Impedance Method)}

\begin{studybox}[Procedure]
\begin{enumerate}[leftmargin=*, label=\textcolor{anthOrange}{\HeadingFont\bfseries\arabic*.}]
  \item Determine the synchronous impedance \(Z_s\) from OCC and short-circuit data.
  \item Calculate \(X_s = \sqrt{Z_s^2 - \Ra^2}\).
  \item Use the phasor equation \(\Ef = \Vt + \Ia(\Ra + jX_s)\).
  \item Obtain \(E_0 = |\Ef|\) and then compute voltage regulation.
\end{enumerate}
\end{studybox}

\begin{formulabox}
\[
  E_0
  =
  \sqrt{\left(\Vt\cos\phi + \Ia\Ra\right)^2
       +
       \left(\Vt\sin\phi \pm \Ia X_s\right)^2}
\]
\(+\) for lagging power factor, \(-\) for leading power factor.
\end{formulabox}

\begin{notebox}
The EMF method usually gives a \textbf{higher} value of regulation and is therefore called a \textbf{pessimistic method}.
\end{notebox}

\subsection{MMF Method (Ampere-Turn Method)}

\begin{studybox}[Idea of the Method]
\begin{itemize}[leftmargin=*]
  \item Convert the required terminal-voltage emf into an equivalent field current from OCC.
  \item Convert the armature reaction and leakage effects into an equivalent field current.
  \item Add these field-current components vectorially according to power factor.
  \item From the resultant field current, read the no-load emf from OCC and compute regulation.
\end{itemize}
\end{studybox}

\begin{notebox}
The MMF method usually gives a \textbf{lower} value of regulation and is therefore regarded as an \textbf{optimistic method}.
\end{notebox}

% ============================================================
\section{Parallel Operation of Alternators and Synchronisation}
% ============================================================

\subsection{Why Alternators are Operated in Parallel}

\begin{itemize}[leftmargin=*]
  \item Better efficiency over varying load.
  \item Greater reliability and continuity of service.
  \item Maintenance and repair can be done without complete shutdown.
  \item Large power demand can be met by combining smaller units.
\end{itemize}

\subsection{Requirements / Desirable Features}

\begin{itemize}[leftmargin=*]
  \item Same output voltage rating.
  \item Speeds such that the machines generate the same frequency.
  \item Same phase sequence and same waveform.
  \item Similar speed-load characteristics of prime movers.
  \item Armature reactance present so current sharing is stable.
\end{itemize}

\subsection{Conditions for Synchronising}

\begin{importantbox}[Conditions to be Satisfied Before Closing the Switch]
\begin{enumerate}[leftmargin=*, label=\textcolor{anthBlue}{\HeadingFont\bfseries\arabic*.}]
  \item Incoming machine terminal voltage must equal bus-bar voltage.
  \item Incoming machine frequency must equal bus-bar frequency.
  \item Phase sequence must be identical.
  \item Phase angle between corresponding voltages must be zero at the instant of switching.
\end{enumerate}
\end{importantbox}

\begin{definebox}[Infinite Bus-Bar]
An \textbf{infinite bus-bar} is an electrical system of practically constant voltage and frequency, unaffected by the connection or disconnection of a single alternator.
\end{definebox}

\subsection{Methods of Synchronisation}

\begin{center}
\begin{tabular}{L{3.7cm} L{5.2cm} L{4.2cm}}
\toprule
\thead{Method} & \thead{Main Observation} & \thead{Switch Closing Instant}\\
\midrule
\trow Dark-lamp method & Lamps become dark when voltages match in phase & Middle of dark period\\
Bright-lamp method & Cross-connected lamps become bright at synchronism reference & Middle of bright period\\
\trow Three-dark-lamp method & Three lamps brighten and darken together if phase sequence is correct & Instant all three are dark together\\
Two-bright-one-dark method & One lamp dark and two bright at synchronism & Lamp \(L_2\) dark, other two bright\\
\trow Synchroscope & Pointer shows relative phase and frequency difference & Pointer at vertical synchronising mark\\
\bottomrule
\end{tabular}
\end{center}

\subsection{Useful Observations from Lamp Methods}

\begin{itemize}[leftmargin=*]
  \item If lamps brighten and dim in sequence, the phase sequence is wrong; interchange any two leads.
  \item In two-bright-one-dark method, a clockwise travelling light wave indicates one speed condition and the opposite direction indicates the other.
  \item The incoming machine is adjusted until lamp changes are very slow, then the switch is closed at the proper instant.
\end{itemize}

% ============================================================
\section{Synchronising Current, Torque and Power}
% ============================================================

\subsection{Synchronising Current}

\begin{itemize}[leftmargin=*]
  \item If two equal emfs are exactly in phase opposition, no circulating current flows.
  \item If one induced emf falls back by a small electrical angle \(\delta\), a resultant emf appears in the local circuit.
\end{itemize}

\begin{formulabox}
For a small angular displacement \(\delta\) (in radians),
\[
  E_R \approx E\delta
\]
and the synchronising current is
\[
  I_{sy} = \frac{E_R}{Z_{cs}} \approx \frac{E\delta}{Z_{cs}}
\]
If circuit resistance is much smaller than reactance,
\[
  I_{sy} \approx \frac{E\delta}{X_{cs}}
\]
\end{formulabox}

\subsection{Role of Synchronising Current}

\begin{itemize}[leftmargin=*]
  \item It acts as a generating current for one machine and a motoring current for the other.
  \item Hence it produces a \textbf{synchronising torque} that accelerates the retarded machine and decelerates the advanced machine.
  \item Any small departure from synchronism therefore produces a restoring tendency.
\end{itemize}

\begin{importantbox}[Synchronising Power]
The power exchanged between two parallel alternators due to synchronising current is called
\textbf{synchronising power}. It is the power-transfer mechanism that helps maintain steady
parallel operation.
\end{importantbox}

% ============================================================
\section{Construction and Principle of Operation of a Synchronous Motor}
% ============================================================

\begin{definebox}[Synchronous Motor]
A \textbf{synchronous motor} is a 3-phase AC motor in which the rotor, excited by DC, runs at the
\textbf{synchronous speed} of the stator rotating magnetic field.
\end{definebox}

\subsection{Construction}

\begin{center}
\begin{tabular}{L{3.0cm} L{9.2cm}}
\toprule
\thead{Part} & \thead{Details}\\
\midrule
\trow Stator / armature & Laminated core, slots on inner periphery, 3-phase winding usually star connected, frame of cast iron or welded steel\\
Rotor / field system & Field poles with field winding supplied by DC, excitation through slip rings and brushes or shaft-mounted exciter\\
\bottomrule
\end{tabular}
\end{center}

\subsection{Working Principle}

\begin{itemize}[leftmargin=*]
  \item The 3-phase stator supply produces a rotating magnetic field at synchronous speed.
  \item The DC-excited rotor behaves like a magnet with fixed polarity.
  \item If the rotor is already near synchronous speed in the same direction as the rotating field, opposite stator and rotor poles attract and lock magnetically.
  \item The rotor then continues to run at synchronous speed with \textbf{zero steady-state slip}.
\end{itemize}

\subsection{Why the Motor is Not Self-Starting}

\begin{studybox}[Reason]
\begin{itemize}[leftmargin=*]
  \item At standstill, the stator poles reverse rapidly every half cycle.
  \item Hence the rotor experiences torque alternately in clockwise and anticlockwise directions.
  \item Due to inertia, the average starting torque is zero.
  \item Therefore a synchronous motor is \textbf{not self-starting} unless a special starting arrangement is provided.
\end{itemize}
\end{studybox}

% ============================================================
\section{Starting Methods of a Synchronous Motor}
% ============================================================

\subsection{Starting by an External DC Motor}

\begin{itemize}[leftmargin=*]
  \item The synchronous motor is mechanically coupled to a DC shunt or compound motor.
  \item The DC motor brings the rotor near synchronous speed.
  \item AC supply is applied to the stator and then DC excitation is applied to the rotor field.
  \item Once magnetic locking takes place, the auxiliary DC motor is disconnected.
\end{itemize}

\subsection{Starting by an AC Motor (Pony Motor)}

\begin{itemize}[leftmargin=*]
  \item A small induction motor is used as a pony motor.
  \item It brings the synchronous motor rotor close to synchronous speed.
  \item After that, the stator AC supply and rotor field excitation are applied for magnetic locking.
\end{itemize}

\subsection{Starting by Damper Winding}

\begin{itemize}[leftmargin=*]
  \item Copper bars embedded in the rotor pole faces and short-circuited at the ends form the \textbf{damper winding}.
  \item At start, the motor behaves like an induction motor.
  \item When it reaches near-synchronous speed, DC excitation is applied and the rotor pulls into step.
\end{itemize}

\begin{importantbox}[Most Practical Method]
Damper winding is the most common practical method because it makes the machine effectively self-starting.
\end{importantbox}

% ============================================================
\section{Equivalent Circuit and Effect of Excitation on Armature Current}
% ============================================================

\subsection{Equivalent Circuit of a Synchronous Motor}

\begin{center}
\begin{circuitikz}[american voltages, scale=0.95]
  \draw
  (0,0) to[sV, l=$V$] (0,3)
  (0,3) to[R, l=$R_a$] (3,3)
        to[L, l=$jX_s$] (6,3)
        to[sV, l_=$E_b$] (6,0)
        -- (0,0);
  \draw[-{Latex[length=2mm]}] (1.2,3.45) -- (5.1,3.45) node[midway, above] {$I_a$};
\end{circuitikz}
\end{center}

\begin{formulabox}
\[
  V = E_b + I_a(R_a + jX_s)
\]
where \(E_b\) is the internal or back emf induced by the rotor field.
\end{formulabox}

\begin{itemize}[leftmargin=*]
  \item The stator is supplied by 3-phase AC.
  \item The rotor is supplied by DC.
  \item Rotor flux induces a back emf \(E_b\) in the stator that opposes the applied voltage.
\end{itemize}

\subsection{Effect of Excitation on Armature Current}

\begin{center}
\begin{tabular}{L{3.5cm} L{4.0cm} L{4.4cm}}
\toprule
\thead{Field Excitation} & \thead{Armature Current} & \thead{Power Factor}\\
\midrule
\trow Under-excited & Higher than minimum & Lagging\\
Normal excitation & Minimum & Approximately unity\\
\trow Over-excited & Increases again & Leading\\
\bottomrule
\end{tabular}
\end{center}

\begin{importantbox}[V-Curve Idea]
At constant load, a plot of armature current versus field current gives a \textbf{V-curve}. The current is minimum at the excitation corresponding to unity power factor.
\end{importantbox}

\begin{notebox}
When over-excited, a synchronous motor draws leading current and can be used for power-factor improvement as a \textbf{synchronous condenser}.
\end{notebox}

% ============================================================
\section{Quick Comparison: Synchronous Generator vs Synchronous Motor}
% ============================================================

\begin{center}
\begin{tabular}{L{3.0cm} L{4.8cm} L{4.8cm}}
\toprule
\thead{Aspect} & \thead{Generator} & \thead{Motor}\\
\midrule
\trow Energy conversion & Mechanical to electrical & Electrical to mechanical\\
Internal emf relation & \(\Ef = \Vt + \Ia(\Ra + jX_s)\) & \(V = E_b + I_a(R_a + jX_s)\)\\
\trow Rotor excitation & DC field required & DC field required\\
Speed & Locked to synchronous speed & Locked to synchronous speed\\
\trow Starting nature & Prime mover starts it naturally & Not self-starting without auxiliary arrangement\\
\bottomrule
\end{tabular}
\end{center}

% ============================================================
\section{Comprehensive Formula Summary}
% ============================================================

\renewcommand{\arraystretch}{1.55}
\begin{longtable}{L{5cm} L{9cm}}
\toprule
\thead{Quantity} & \thead{Formula / Key Relation}\\
\midrule
\endfirsthead
\multicolumn{2}{c}{\small\itshape (continued)}\\
\toprule
\thead{Quantity} & \thead{Formula / Key Relation}\\
\midrule
\endhead
\bottomrule
\endlastfoot

\trow\textbf{Synchronous speed} & \(N_s = \dfrac{120f}{P}\)\\
\textbf{Generated frequency} & \(f = \dfrac{PN_s}{120}\)\\
\trow\textbf{Pitch factor} & \(\kp = \cos\dfrac{\alpha}{2}\)\\
\textbf{Distribution factor} & \(\kd = \dfrac{\sin(m\beta/2)}{m\sin(\beta/2)}\)\\
\trow\textbf{Winding factor} & \(\kw = \kp\kd\)\\
\textbf{Alternator emf per phase} & \(\Eph = 4.44\,\kw\,f\,\Phim\,T\)\\
\trow\textbf{Synchronous reactance} & \(X_s = X_l + X_a\)\\
\textbf{Synchronous impedance} & \(Z_s = \sqrt{R_a^2 + X_s^2}\)\\
\trow\textbf{Generator phasor equation} & \(\Ef = \Vt + \Ia(\Ra + jX_s)\)\\
\textbf{Voltage regulation} & \(\%VR = \dfrac{E_0-\Vt}{\Vt}\times 100\)\\
\trow\textbf{EMF method} & \(E_0 = \sqrt{(\Vt\cos\phi + \Ia \Ra)^2 + (\Vt\sin\phi \pm \Ia X_s)^2}\)\\
\textbf{Armature reaction at UPF} & Cross-magnetising\\
\trow\textbf{Armature reaction at ZPF lagging} & Demagnetising\\
\textbf{Armature reaction at ZPF leading} & Magnetising\\
\trow\textbf{Resultant emf for small angular displacement} & \(E_R \approx E\delta\)\\
\textbf{Synchronising current} & \(I_{sy} = \dfrac{E_R}{Z_{cs}} \approx \dfrac{E\delta}{Z_{cs}}\)\\
\trow\textbf{Motor equivalent equation} & \(V = E_b + I_a(R_a + jX_s)\)\\
\textbf{Unity-pf excitation point} & Minimum armature current on V-curve\\
\end{longtable}
\renewcommand{\arraystretch}{1.0}

% ============================================================
\section{Exam-Focused Revision Points}
% ============================================================

\begin{importantbox}[High-Yield Points]
\begin{enumerate}[leftmargin=*, label=\textcolor{anthOrange}{\HeadingFont\bfseries\arabic*.}]
  \item In alternators, the \textbf{field} is on the rotor and the \textbf{armature} is on the stator.
  \item Salient-pole rotors are for \textbf{low-speed} machines; cylindrical rotors are for \textbf{high-speed} turbo alternators.
  \item Generated emf per phase depends on \(\kw\), frequency, flux per pole and turns per phase.
  \item At unity pf, armature reaction is mainly \textbf{cross-magnetising}; at lagging pf it is demagnetising; at leading pf it is magnetising.
  \item EMF method is \textbf{pessimistic}; MMF method is \textbf{optimistic}.
  \item Synchronising requires equality of \textbf{voltage, frequency, phase sequence and phase angle}.
  \item A synchronous motor is \textbf{not self-starting}; damper winding solves this practically.
  \item Over-excited synchronous motors draw \textbf{leading current} and help improve power factor.
\end{enumerate}
\end{importantbox}

% ============================================================
\section{Original Slide Appendix}
% ============================================================

\addcontentsline{toc}{section}{Original Slide Appendix}
\begin{studybox}[Why the Appendix is Included]
The notes above are rewritten for study. The original source slides are embedded below so that
no diagram-only or slide-specific content is lost from the final single PDF.
\end{studybox}

\includepdf[pages=-,fitpaper=true,pagecommand={\thispagestyle{plain}}]{build/part1.pdf}
\includepdf[pages=-,fitpaper=true,pagecommand={\thispagestyle{plain}}]{build/part2.pdf}
\includepdf[pages=-,fitpaper=true,pagecommand={\thispagestyle{plain}}]{build/part3.pdf}

% ── back cover ───────────────────────────────────────────────
\clearpage
\thispagestyle{empty}
\vfill
\begin{center}
\begin{tikzpicture}
  \fill[anthDark] (0,0) rectangle (\textwidth, 0.52cm);
  \fill[anthOrange] (0,0.52cm) rectangle (\textwidth, 0.68cm);
  \node[anthLight, font=\HeadingFont\bfseries\small] at (0.5\textwidth, 0.26cm)
    {End of Module 3 — Synchronous Machines\quad|\quad PCEOT402 Electrical Machines};
\end{tikzpicture}
\end{center}

\end{document}
